\section{Proofs from Section \ref{subsec:finite-dim}}
\label{sec:proof-finite-dim}

Throughout this section,
Assumption~\ref{assumption:finite-dim-loss-penalty} holds.
We first prove the intermediate propositions in
Section~\ref{subsec:proof-intermediate-res-finite-dim},
then establish global descent in
Section~\ref{subsec:global-descent}, finite branch identification
in Section~\ref{subsec:branch-identification}, and convergence
in Section~\ref{subsec:convergence-within-branch}.

\subsection{Proof of intermediate results}
\label{subsec:proof-intermediate-res-finite-dim}

\begin{proof}[Proof of Proposition~\ref{prop:ST-potential}]
Direct integration gives
\begin{equation}
    \Psi_\kappa(r)
    =
    \begin{cases}
        \ell(-1,r)-\ell(-1,\kappa_-),
            & r\leq\kappa_-,\\
        0,
            & \kappa_-<r<\kappa_+,\\
        \ell(+1,r)-\ell(+1,\kappa_+),
            & r\geq\kappa_+.
    \end{cases}
    \label{eq:psi-kappa-formula}
\end{equation}
Nonnegativity and monotonicity of the loss give
\begin{align}
    -C_\kappa\leq\Psi_\kappa\leq0,
    \qquad
    C_\kappa:=
    \max\{\ell(-1,\kappa_-),\ell(+1,\kappa_+)\}. \label{eq:c-kappa}
\end{align}
Thus $\Psi_\kappa$ is bounded and continuous, and is twice
continuously differentiable away from the thresholds.
Convexity and the derivative signs give
\[
    |\psi_\kappa(r)|
    \leq
    \max\{\ell'(-1,\kappa_-),-\ell'(+1,\kappa_+)\}
    <\infty,
\]
so its integral $\Psi_\kappa$ is globally Lipschitz.

Since the supervised loss and penalty are $C^2$, it follows that
$V$ is continuous and $C^2$ outside $\mathcal S_\kappa$.
There, the chain rule gives
\[
    \nabla V(\theta)
    =
    \nabla L_{\mathcal L}(\theta)
    +
    \frac{\pi}{M}\sum_{j\in I_{\mathcal U}}
        \psi_\kappa(r_j(\theta))\breve X_j
    +
    \frac{\lambda}{d}
        \begin{bmatrix}0\\\nabla J(w)\end{bmatrix}
    =
    G(\theta).
\]
Each threshold equation defines a proper affine hyperplane,
since the coefficient of $b$ in $r_j(\theta)$ is one, and
their finite union has Lebesgue measure zero.
\end{proof}

\begin{proof}[Proof of Proposition~\ref{prop:configuration-branch-formulation}]
By the selection convention, $\theta\in\Theta_{\mathfrak c}$
if and only if, for every $j\in I_{\mathcal U}$,
\[
    \begin{cases}
        r_j(\theta)\leq\kappa_-,
            & \mathfrak c_j=-1,\\
        \kappa_-<r_j(\theta)<\kappa_+,
            & \mathfrak c_j=0,\\
        r_j(\theta)\geq\kappa_+,
            & \mathfrak c_j=+1.
    \end{cases}
\]
Hence each branch is convex.
For a nonempty branch, its closure is obtained by replacing
the strict inequalities by non-strict ones.
Indeed, fix a point in the branch and interpolate it with
any point satisfying these weak inequalities.
Every interpolant with positive weight on the branch point
satisfies the strict inequalities and remains in the branch.
The weak-inequality set is closed, proving the claimed
closure characterization. It is therefore polyhedral.
The partition property follows from the definition of
$\mathfrak C$, whose range contains at most $3^M$ configurations.

The explicit formula \eqref{eq:psi-kappa-formula} gives
$V=L_{\mathfrak c}+K_{\mathfrak c}$ on $\Theta_{\mathfrak c}$,
where
\[
    K_{\mathfrak c}
    :=
    -\frac{\pi}{M}
    \left[
        \#\{j\in I_{\mathcal U}:\mathfrak c_j=-1\}
            \ell(-1,\kappa_-)
        +
        \#\{j\in I_{\mathcal U}:\mathfrak c_j=+1\}
            \ell(+1,\kappa_+)
    \right].
\]
Continuity extends the identity to
$\overline{\Theta_{\mathfrak c}}$.
Finally, for every $\vartheta\in\Theta_{\mathfrak c}$,
\[
    L_{\mathrm{ST}}(\theta;\vartheta)
    =
    L_{\mathfrak c}(\theta)
    \qquad\text{for every }\theta.
\]
Differentiating with respect to $\theta$ and then setting
$\vartheta=\theta\in\Theta_{\mathfrak c}$ gives
$G(\theta)=\nabla L_{\mathfrak c}(\theta)$, including at thresholds
assigned to that branch.
\end{proof}

\begin{proof}[Proof of Proposition~\ref{prop:branch-convexity-coercivity}]
By Assumption~\ref{assumption:finite-dim-loss-penalty},
$\ell''(y,r)>0$ and $j''(u)>0$.
Since $\pi\geq0$, for every nonzero direction
$(a,u)\in\mathbb R\times\mathbb R^d$,
\[
    \begin{bmatrix}a\\u\end{bmatrix}^{\!\top}
    \nabla^2L_{\mathfrak c}(\theta)
    \begin{bmatrix}a\\u\end{bmatrix}
    \geq
    \frac1N\sum_{i\in I_{\mathcal L}}
        \ell''(Y_i,r_i(\theta))
        \left(a+\frac{X_i^\top u}{\sqrt d}\right)^2
    +
    \frac{\lambda}{d}\sum_{k=1}^d j''(w_k)u_k^2
    >0.
\]
Indeed, the penalty term is positive when $u\neq0$, while the
labeled term is positive when $u=0$ and $a\neq0$.
Thus every branch Hessian is positive definite, and every
$L_{\mathfrak c}$ is strictly convex.
Moreover, $\ell''\leq L_\ell$, $j''\leq L_j$, and $\pi\geq0$
imply
\[
    \nabla^2L_{\mathfrak c}(\theta)\preceq H_d.
\]
Indeed, the unlabeled Hessian contribution is bounded by
enlarging its sum to all unlabeled observations.

The sublevel sets of the supervised objective
$L_{\mathcal L}(\theta)+(\lambda/d)J(w)$ bound $w$
through the nonnegative, coercive penalty.
Convexity and the strict derivative signs imply
\[
    \ell(-1,r)\longrightarrow+\infty
        \quad\text{as }r\longrightarrow+\infty,
    \qquad
    \ell(+1,r)\longrightarrow+\infty
        \quad\text{as }r\longrightarrow-\infty.
\]
With $w$ bounded, a negative labeled observation therefore prevents
$b\to+\infty$, and a positive labeled observation prevents
$b\to-\infty$.
The supervised objective is consequently coercive.
Every $L_{\mathfrak c}$ is bounded below by this objective,
while $V$ differs from it by a bounded term.
Consequently, both are coercive.
Continuity gives existence of each branch minimizer, and
strict convexity gives uniqueness.
\end{proof}

\subsection{Proof of Theorem~\ref{thm:st-convergence}}

\subsubsection{Global descent}
\label{subsec:global-descent}

\begin{lemma}[Global descent and boundedness]
\label{lem:st-global-descent-boundedness}
For $\pi,\lambda\geq0$, the potential satisfies
\begin{equation}
    V(\theta+h)-V(\theta)
    \leq
    G(\theta)^\top h+\frac12h^\top H_dh,
    \qquad \theta,h\in\mathbb R^{d+1}.
    \label{eq:st-global-quadratic-upper-bound}
\end{equation}
Under \eqref{eq:st-descent-step-size}, set
$c_d:=\eta(1-\eta\Lambda_d/2)>0$.
Then the iterates satisfy
\begin{equation}
    V(\theta^{t+1})-V(\theta^t)
    \leq
    -c_d\|P_d^{1/2}G(\theta^t)\|^2,
    \qquad t\geq1,
    \label{eq:st-global-descent}
\end{equation}
and $G(\theta^t)\to0$.
If, in addition, $\lambda>0$ and both classes occur in the
labeled sample, then the trajectory is bounded.
\end{lemma}
\begin{proof}
Since $\ell''(y,r)\leq L_\ell$, for every $r,\delta\in\mathbb R$,
\[
    \ell(y,r+\delta)
    \leq
    \ell(y,r)+\ell'(y,r)\delta+\frac{L_\ell}{2}\delta^2.
\]
On each smooth interval,
$0\leq\psi_\kappa'(u)\leq L_\ell$.
The signs
$\ell'(-1,\kappa_-)>0$ and $\ell'(+1,\kappa_+)<0$
give a downward jump at each threshold.
With the prescribed values at the thresholds,
$u\mapsto\psi_\kappa(u)-L_\ell u$ is nonincreasing.
Therefore
\[
    (\psi_\kappa(v)-\psi_\kappa(u))(v-u)
    \leq L_\ell(v-u)^2,
    \qquad u,v\in\mathbb R.
\]
The definition of $\Psi_\kappa$ gives
\begin{align*}
    \Psi_\kappa(r+\delta)-\Psi_\kappa(r)
        -\psi_\kappa(r)\delta
    =
    \int_0^1[\psi_\kappa(r+s\delta)-\psi_\kappa(r)]
        \delta\,ds
    \leq
    \int_0^1L_\ell s\delta^2\,ds
    =
    \frac{L_\ell}{2}\delta^2.
\end{align*}
For $s>0$, the integrand bound follows by applying the preceding
inequality with $u=r$ and $v=r+s\delta$, then dividing by $s$.
The argument holds for either sign of $\delta$ and across thresholds.

Apply these scalar bounds with $\delta=\breve X_i^\top h$.
Writing $h=(h_b,h_w)\in\mathbb R\times\mathbb R^d$, also use
\[
    J(w+h_w)-J(w)
    \leq
    \nabla J(w)^\top h_w+\frac{L_j}{2}\|h_w\|^2.
\]
Summing, with $\pi,\lambda\geq0$, proves
\eqref{eq:st-global-quadratic-upper-bound}.
Substituting $h=-\eta P_dG(\theta^t)$ and using the definition of
$\Lambda_d$ yields \eqref{eq:st-global-descent}.
In particular, the potential decreases strictly at every
nonstationary iteration, including iterations that cross
branch boundaries.

The potential is bounded below because $\Psi_\kappa$ is bounded
and the loss and penalty are nonnegative.
Summing the descent inequality therefore gives
\[
    c_d\sum_{t=1}^T\|P_d^{1/2}G(\theta^t)\|^2
    \leq V(\theta^1)-V(\theta^{T+1})
    \leq V(\theta^1)-\inf_\theta V(\theta)<\infty.
\]
Consequently, writing
$\Delta\theta^t:=\theta^{t+1}-\theta^t$,
\begin{equation}
    \sum_{t=1}^\infty\|P_d^{1/2}G(\theta^t)\|^2<\infty,
    \qquad
    \sum_{t=1}^\infty\|\Delta\theta^t\|_{P_d^{-1}}^2<\infty,
    \label{eq:st-square-summable-updates}
\end{equation}
where the second sum is $\eta^2$ times the first.
For fixed $d$, invertibility of $P_d$ and the update identity imply
\begin{equation}
    G(\theta^t)\xrightarrow[t\to\infty]{}0,
    \qquad
    \Delta\theta^t=-\eta P_dG(\theta^t)
        \xrightarrow[t\to\infty]{}0.
    \label{eq:parameter-difference-converges-to-zero}
\end{equation}
If $\lambda>0$ and both classes occur in the labeled sample,
Proposition~\ref{prop:branch-convexity-coercivity} gives coercivity
of $V$. By descent, the iterates from $t=1$ onward remain in
the compact sublevel set
\[
    \{\theta\in\mathbb R^{d+1}:V(\theta)\leq V(\theta^1)\}.
\]
Thus the trajectory is bounded.
\end{proof}

\subsubsection{Branch identification}
\label{subsec:branch-identification}

\begin{lemma}[Finite branch identification]
\label{lem:st-finite-identification}
Suppose $\pi>0$.
If a trajectory of \eqref{eq:ST-recursion-G} is bounded and satisfies
$G(\theta^t)\to0$, then there are a finite time
$T_{\mathrm{id}}\geq1$ and a nonempty branch
$\Theta_{\mathfrak c}$ such that
\begin{equation}
    \mathfrak C(\theta^t)=\mathfrak c,
    \qquad t\geq T_{\mathrm{id}}.
    \label{eq:st-finite-branch-identification}
\end{equation}
\end{lemma}

\begin{proof}
We first show that gradients of distinct branches cannot coincide
at a point belonging to both branch closures.
Let $\Theta_{\mathfrak c}$ and $\Theta_{\mathfrak d}$ be nonempty,
distinct branches, and let
\[
    \theta^\star\in
        \overline{\Theta_{\mathfrak c}}
        \cap\overline{\Theta_{\mathfrak d}},
    \qquad
    D_{\mathfrak c,\mathfrak d}
    :=
    \{j\in I_{\mathcal U}:\mathfrak c_j\neq\mathfrak d_j\}.
\]
Since $\kappa_-<\kappa_+$, the closures of the negative and positive
selection regions are disjoint.
Thus, for each $j\in D_{\mathfrak c,\mathfrak d}$, exactly one of
$\mathfrak c_j,\mathfrak d_j$ is zero, and $r_j(\theta^\star)$
equals the threshold corresponding to the nonzero state.

Set
\[
    \alpha_j
    :=
    \left|
        \ell'\bigl(
            \operatorname{sign}(r_j(\theta^\star)),
            r_j(\theta^\star)
        \bigr)
    \right|
    >0,
    \qquad j\in D_{\mathfrak c,\mathfrak d}.
\]
The supervised and penalty gradients, and the unlabeled terms
with equal states, cancel. Since the nonzero state has the same
sign as $r_j(\theta^\star)$, the derivative signs give
\[
    \nabla L_{\mathfrak c}(\theta^\star)
    -
    \nabla L_{\mathfrak d}(\theta^\star)
    =
    -\frac{\pi}{M}
    \sum_{j\in D_{\mathfrak c,\mathfrak d}}
        (\mathfrak c_j-\mathfrak d_j)\alpha_j\breve X_j.
\]
Choose $\theta_{\mathfrak c}\in\Theta_{\mathfrak c}$ and
$\theta_{\mathfrak d}\in\Theta_{\mathfrak d}$.
The branch inequalities imply
\[
    (\mathfrak c_j-\mathfrak d_j)
    \bigl(
        r_j(\theta_{\mathfrak c})
        -
        r_j(\theta_{\mathfrak d})
    \bigr)
    =
    |r_j(\theta_{\mathfrak c})-r_j(\theta_{\mathfrak d})|
    >0,
    \qquad j\in D_{\mathfrak c,\mathfrak d},
\]
because the selected score lies at or beyond its threshold,
whereas the unselected score lies strictly inside the threshold
interval.
Taking the inner product with
$\theta_{\mathfrak c}-\theta_{\mathfrak d}$ therefore gives
\begin{equation}
\begin{aligned}
    \bigl(
        \nabla L_{\mathfrak c}(\theta^\star)
        -
        \nabla L_{\mathfrak d}(\theta^\star)
    \bigr)^\top
    (\theta_{\mathfrak c}-\theta_{\mathfrak d})
    =
    -\frac{\pi}{M}
    \sum_{j\in D_{\mathfrak c,\mathfrak d}}
        \alpha_j
        |r_j(\theta_{\mathfrak c})-r_j(\theta_{\mathfrak d})|
    <0.
\end{aligned}
    \label{eq:st-branch-gradient-separation}
\end{equation}
Strictness follows from $\pi>0$,
$D_{\mathfrak c,\mathfrak d}\neq\varnothing$, and $\alpha_j>0$.

If configurations changed infinitely often, finiteness of the
configuration set would give distinct configurations
$\mathfrak c,\mathfrak d$ and times $t_k\to\infty$ such that
\[
    \theta^{t_k}\in\Theta_{\mathfrak c},
    \qquad
    \theta^{t_k+1}\in\Theta_{\mathfrak d}.
\]
By boundedness, after passing to a subsequence,
$\theta^{t_k}\to\theta^\star$.
The update identity and $G(\theta^t)\to0$ imply
$\theta^{t+1}-\theta^t\to0$, so
$\theta^{t_k+1}\to\theta^\star$ as well.
Hence $\theta^\star$ belongs to both branch closures.
By Proposition~\ref{prop:configuration-branch-formulation}
and continuity of the branch gradients,
\[
    \nabla L_{\mathfrak c}(\theta^\star)
    =
    \lim_{k\to\infty}G(\theta^{t_k})=0,
    \qquad
    \nabla L_{\mathfrak d}(\theta^\star)
    =
    \lim_{k\to\infty}G(\theta^{t_k+1})=0.
\]
This contradicts \eqref{eq:st-branch-gradient-separation}.
There are therefore only finitely many configuration changes,
proving \eqref{eq:st-finite-branch-identification}.
\end{proof}

\subsubsection{Convergence within the identified branch}
\label{subsec:convergence-within-branch}

\begin{proof}[Proof of Theorem~\ref{thm:st-convergence}]
Lemma~\ref{lem:st-global-descent-boundedness} gives boundedness and
$G(\theta^t)\to0$.
Lemma~\ref{lem:st-finite-identification} then gives
$T_{\mathrm{id}}<\infty$.
For the identified branch $\Theta_{\mathfrak c}$,
Proposition~\ref{prop:configuration-branch-formulation} yields
\[
    \theta^{t+1}
    =
    \theta^t-\eta P_d\nabla L_{\mathfrak c}(\theta^t),
    \qquad t\geq T_{\mathrm{id}}.
\]
Boundedness guarantees at least one accumulation point.
Every such point belongs to $\overline{\Theta_{\mathfrak c}}$
and is a zero of $\nabla L_{\mathfrak c}$, by continuity of the
branch gradient and $G(\theta^t)\to0$.
By Proposition~\ref{prop:branch-convexity-coercivity},
every accumulation point therefore equals the unique global
minimizer $\theta_{\mathfrak c}^\star$.
Consequently,
\begin{equation}
    \theta^t\longrightarrow\theta_{\mathfrak c}^\star,
    \qquad
    \nabla L_{\mathfrak c}(\theta_{\mathfrak c}^\star)=0,
    \qquad
    \theta_{\mathfrak c}^\star
        \in\overline{\Theta_{\mathfrak c}}.
    \label{eq:st-branch-limit}
\end{equation}
For the geometric rate, work in the coordinates
$x^t=P_d^{-1/2}\theta^t$ and define
\[
    T_{\mathfrak c}(x)
    :=
    x-\eta P_d^{1/2}
        \nabla L_{\mathfrak c}(P_d^{1/2}x),
    \qquad
    x_{\mathfrak c}^\star
    :=
    P_d^{-1/2}\theta_{\mathfrak c}^\star.
\]
Then $x^{t+1}=T_{\mathfrak c}(x^t)$ for
$t\geq T_{\mathrm{id}}$, and
$T_{\mathfrak c}(x_{\mathfrak c}^\star)
=x_{\mathfrak c}^\star$.
Let
\[
    K
    :=
    \overline{\operatorname{conv}}
    \left(
        \{x^t:t\geq T_{\mathrm{id}}\}
        \cup\{x_{\mathfrak c}^\star\}
    \right).
\]
Boundedness of the trajectory implies that $K$ is compact.
For $x\in K$, write
\[
    A(x)
    :=
    P_d^{1/2}
        \nabla^2L_{\mathfrak c}(P_d^{1/2}x)
    P_d^{1/2}.
\]
The branch Hessian is continuous and positive definite
everywhere. Compactness therefore gives
\[
    m_K
    :=
    \min_{x\in K}\lambda_{\min}(A(x))
    >0.
\]
Moreover, the upper Hessian bound and the definition of
$\Lambda_d$ give
\[
    m_K I_{d+1}
    \preceq A(x)
    \preceq \Lambda_d I_{d+1},
    \qquad x\in K.
\]
Since
\(
    DT_{\mathfrak c}(x)
    =
    I_{d+1}-\eta A(x)
\),
the step-size condition \eqref{eq:st-descent-step-size}
implies
\[
    \sup_{x\in K}
        \|DT_{\mathfrak c}(x)\|_{\mathrm{op}}
    \leq
    \max\{|1-\eta m_K|,\ |1-\eta\Lambda_d|\}
    <1.
\]
Choose $q\in(0,1)$ bounding the left-hand side.
Convexity of $K$ and the integral form of the mean-value
formula yield
\[
    T_{\mathfrak c}(x)-T_{\mathfrak c}(y)
    =
    \int_0^1
        DT_{\mathfrak c}(y+s(x-y))(x-y)\,ds,
    \qquad x,y\in K,
\]
and hence
\begin{equation}
    \|T_{\mathfrak c}(x)-T_{\mathfrak c}(y)\|
    \leq q\|x-y\|,
    \qquad x,y\in K.
    \label{eq:st-branch-contraction}
\end{equation}
The set $K$ need not lie inside the identified branch:
the smooth map $T_{\mathfrak c}$ is defined everywhere,
and the actual iterates already remain in that branch.
Applying \eqref{eq:st-branch-contraction} with
$x=x^t$ and $y=x_{\mathfrak c}^\star$ gives
\[
    \|x^{t+1}-x_{\mathfrak c}^\star\|
    \leq q\|x^t-x_{\mathfrak c}^\star\|,
    \qquad t\geq T_{\mathrm{id}}.
\]
Iteration and the identity
$\|P_d^{-1/2}v\|=\|v\|_{P_d^{-1}}$ now give
\[
    \|\theta^t-\theta_{\mathfrak c}^\star\|_{P_d^{-1}}
    \leq
    q^{\,t-T_{\mathrm{id}}}
    \|\theta^{T_{\mathrm{id}}}
        -\theta_{\mathfrak c}^\star\|_{P_d^{-1}},
    \qquad t\geq T_{\mathrm{id}}.
\]
Finally,
\[
    \|\theta^{t+1}-\theta^t\|_{P_d^{-1}}
    \leq
    (1+q)q^{\,t-T_{\mathrm{id}}}
    \|\theta^{T_{\mathrm{id}}}
        -\theta_{\mathfrak c}^\star\|_{P_d^{-1}}.
\]
The geometric series and the finite initial segment show
that the trajectory has finite length.
\end{proof}
